% TOP_PROBLEM: {"id": 3308, "title": "Coloring random subgraphs", "category_id": 3, "category": {"id": 3, "name": "graph_theory", "display_name": "Graph Theory", "description": "Problems involving graphs, networks, and their properties.", "slug": "graph-theory", "order_index": 3, "created_at": "2026-07-31T15:26:25.671Z"}, "status": "open", "problem_number": "OPG-827", "set_id": 12, "statusQualification": "The logarithmic expression is restricted to graphs with chromatic number at least two."}
% FIRST_POSTED: 2026-10-01
% ENTRY_KIND: statement_only
% UnsolvedMath ID 3308; source catalogue code OPG-827
% !TeX program = lualatex
% Definition and sources only; this file is not an LLM solution attempt.
\documentclass[11pt,a4paper]{article}
\usepackage[margin=25mm]{geometry}
\usepackage{fontspec}
\setmainfont{lmroman10-regular.otf}[BoldFont=lmroman10-bold.otf,
 ItalicFont=lmroman10-italic.otf,BoldItalicFont=lmroman10-bolditalic.otf]
\usepackage{amsmath,amssymb,mathtools}
\usepackage{unicode-math}
\setmathfont{latinmodern-math.otf}
\AtBeginDocument{\let\setminus\smallsetminus}
\usepackage{xurl,hyperref}
\hypersetup{colorlinks=true,urlcolor=blue}
\setcounter{secnumdepth}{0}
\setlength{\parindent}{0pt}
\setlength{\parskip}{5pt}
\setlength{\emergencystretch}{3em}
\begin{document}
\section*{Coloring random subgraphs}\label{3308}
{\small UnsolvedMath ID 3308 · OPG-827 · Graph Theory · OpenGarden}
\subsection{Definitions and mathematical statement}
Let $G=(V,E)$ be a finite simple graph containing at least one edge,
and let $k=\chi(G)\ge2$ be its chromatic number. A proper colouring
assigns different colours to adjacent vertices. Form the random
spanning subgraph $G_{1/2}=(V,E')$ by retaining each edge of $G$
independently with probability $1/2$. The vertex set is unchanged,
even when retained edges leave some vertices isolated.

The problem asks whether there exists an absolute constant $c>0$
such that every such graph satisfies
\[
                   \mathbb E\bigl[\chi(G_{1/2})\bigr]
                      >c\,\frac{\chi(G)}{\log\chi(G)}.
\]
Expectation is over all $2^{|E|}$ equally likely edge subsets, and
$\log$ means the natural logarithm. Changing the logarithm base only
rescales the unknown constant. The restriction $k\ge2$ explicitly
avoids the undefined denominator for a nonempty edgeless graph.

Equivalently, the question is whether
$\inf_G\mathbb E[\chi(G_{1/2})]\log\chi(G)/\chi(G)>0$, with the
infimum over finite graphs having an edge. The strict inequality can
be replaced by a non-strict one after decreasing the constant.

The same $c$ must work for all orders and all graph structures. A
bound involving $\log|V|$ instead of $\log\chi(G)$ is a different
claim, since these quantities can be arbitrarily far apart. The
random operation deletes edges, not vertices, and no conditioning
on connectedness or on the retained number of edges is imposed.
\subsection{Short English statement}
Does independently keeping half the edges preserve an expected chromatic number at least a constant times the original chromatic number divided by its logarithm?
\subsection{Sources}
\begin{itemize}
\item[S1] ULAM.AI and the UnsolvedMath Contributors, supplied problems.json, record 3308 (OPG-827), OpenGarden collection. Catalogue identity and source transcription; CC BY 4.0.
\url{https://huggingface.co/datasets/ulamai/UnsolvedMath}.
\item[S2] Open Problem Garden, Coloring random subgraphs. Original problem and discussion; the mathematical formulation below is expanded from the supplied source record.
\url{http://www.openproblemgarden.org/op/coloring_random_subgraphs}.
\end{itemize}
\end{document}
